\label{sec:longitudinal}

\subsection{Durability of Competitive Districts Under Vote-Share Drift}

A competitive district at a given moment may not remain competitive as
national partisan tides shift.  A district that is D+8 in a year when
Democrats win the presidential popular vote by 3 points could become
uncompetitive if the national environment shifts by 5 points toward
Republicans.  We assess durability by applying uniform partisan shifts of
$\pm 1$ through $\pm 4$ percentage points to the 2020 presidential baseline.

Formally, the shifted Democratic vote share in district $i$ under shift $s$ is:
\begin{equation}
  D_i(s) = D_i + s, \quad s \in \{-4, -3, -2, -1, +1, +2, +3, +4\}
\end{equation}
We count the number of algorithmic districts that remain competitive
($|D_i(s) - 0.5| < 0.10$) under each shift.

\subsection{Results: Longitudinal Stability}

Table~\ref{tab:longitudinal} presents the count of competitive algorithmic
districts as national vote shares drift.

\begin{table}[h]
\centering
\caption{Competitive Algorithmic Districts Under Partisan Drift}
\label{tab:longitudinal}
\begin{tabular}{lcc}
\toprule
National shift & Competitive districts & Swing districts \\
\midrule
$-4$ pts (R wave) & 71 & 42 \\
$-3$ pts & 75 & 47 \\
$-2$ pts & 79 & 51 \\
$-1$ pt & 82 & 54 \\
Baseline (2020) & 85 & 57 \\
$+1$ pt (D wave) & 83 & 55 \\
$+2$ pts & 80 & 52 \\
$+3$ pts & 77 & 49 \\
$+4$ pts & 73 & 44 \\
\bottomrule
\end{tabular}
\end{table}

The count of competitive districts declines modestly under large partisan
shifts but remains substantially higher than the enacted baseline (65) even
under a 4-point wave in either direction.  This indicates that competitive
algorithmic districts are not clustered just inside the 10-point threshold
in a way that would make them sensitive to small shifts.  The distribution
of margins within competitive algorithmic districts---median 4.1 points, with
many at 2--3 points---creates a broad competitive core that is stable across
moderate swings.

\subsection{Durably Competitive Districts}

A district is \emph{durably competitive} if it remains competitive under
every shift from $-4$ to $+4$ points.  Under this strict criterion, 61 of
the 85 competitive algorithmic districts (72\%) are durably competitive.
This compares favorably with enacted maps, where only 44 of the 65 competitive
enacted districts (68\%) are durably competitive.

The 24 competitive algorithmic districts that are not durably competitive
are clustered near the 10-point boundary---they become uncompetitive under a
4-point shift in one direction.  These borderline cases represent districts
where the underlying geography is somewhat homogeneous and which would require
an unusually moderate electorate to remain contested.

\subsection{Projection to Future Census Cycles}

\citet{dellaLibera2026stability} shows that algorithmic district boundaries
have strong temporal stability across census decades.  For competitive
districts specifically, this means that a district that is competitive in the
2020 cycle is likely to have similar geographic properties in the 2030 cycle,
with boundary changes driven primarily by population shifts within the
district's territory rather than by map redesign.  The durability result
therefore extends beyond a single decade: competitive algorithmic districts
represent durable geographic features of the electoral landscape, not
artifacts of a specific partisan moment.
